\section{Evaluation}
\label{sec:evaluation}

\subsection{Experimental Setup}

We apply GeoSection to all 50 states using 2020 Census data, TIGER tract
geometries, and 2020 presidential election results
(precinct-to-tract interpolation via VEST/Fekrazad \citep{fekrazad2021}).
District counts follow the 2020 apportionment (435 total seats).
All runs use 50 seeds per ratio and $\lambda = 0$ (Phase 1 only).
Population balance tolerance is 0.5\,\% ($\delta = 0.005$), matching the
congressional standard established in \emph{Karcher v.\ Daggett} (1983).

The reference implementation is the \texttt{redist} Rust binary using
METIS 5.1 (vendored via \texttt{metis-rs} 0.2 FFI crate).
All 45 multi-district states converge within the 0.5\,\% balance tolerance.
Five single-district states (AK, DE, ND, SD, VT, WY) receive trivial
single-district assignments and are excluded from the ratio analysis.

The MEC baseline (B.7) uses the same pipeline with standard
$\lfloor k/2 \rfloor:\lceil k/2 \rceil$ bisection and 50 seeds per level.
All partisan analysis uses 2020 presidential two-party vote share.

\subsection{50-State Natural Ratio Results}
\label{sec:sweep}

Table~\ref{tab:natural-ratios} reports the GeoSection natural ratio for
every multi-district state, alongside the B.7 MEC baseline seat counts.

\begin{table}[h]
\centering
\caption{GeoSection natural ratios, 50-state sweep.
  ``Equal/near-equal'' = within $20\,\%$ of $k/2:k/2$.
  EC = top-level edge-cut at natural ratio. D\,seats = Democratic congressional seats.
  Columns: State, Seats ($k$), Natural ratio, EC (km), D seats (GeoSection), D seats (MEC), D-vote \%.}
\label{tab:natural-ratios}
\small
\begin{tabular}{llrlrrr}
\toprule
State & $k$ & Natural ratio & EC (km) & D seats & MEC D & D vote \% \\
\midrule
\multicolumn{7}{l}{\textit{Equal or near-equal splits (normalisation dominant)}} \\
IA  &  4 & 2:2 &  962 & 1 & 1 & 45.8 \\
MS  &  4 & 1:3 &  903 & 1 & 1 & 41.6 \\
MD  &  8 & 4:4 &  836 & 6 & 6 & 67.0 \\
MN  &  8 & 1:7 & 1675 & 4 & 3 & 53.6 \\
MA  &  9 & 4:5 &  682 & 9 & 9 & 67.1 \\
SC  &  7 & 3:4 & 1445 & 1 & 3 & 44.1 \\
NJ  & 12 & 5:7 &  753 & 9 & 10 & 58.1 \\
NC  & 14 & 6:8 & 2510 & 5 & 5 & 49.3 \\
NY  & 26 & 11:15 & 2295 & 21 & 21 & 61.7 \\
\midrule
\multicolumn{7}{l}{\textit{Asymmetric splits (genuine urban compactness confirmed)}} \\
MT  &  2 & 1:1 &  733 & 0 & 0 & 41.6 \\
NE  &  3 & 1:2 &  495 & 1 & 1 & 40.2 \\
NM  &  3 & 1:2 & 1148 & 2 & 2 & 55.5 \\
KS  &  4 & 1:3 &  878 & 1 & 1 & 42.5 \\
NV  &  4 & 1:3 &  753 & 2 & 2 & 51.2 \\
UT  &  4 & 1:3 &  960 & 1 & 1 & 39.3 \\
AR  &  4 & 1:3 & 1089 & 0 & 0 & 35.8 \\
CT  &  5 & 1:4 &  432 & 5 & 5 & 60.2 \\
OK  &  5 & 1:4 & 1187 & 1 & 0 & 33.1 \\
KY  &  6 & 1:5 & 1025 & 1 & 1 & 36.8 \\
OR  &  6 & 1:5 & 1262 & 4 & 4 & 58.3 \\
LA  &  6 & 1:5 & 1437 & 1 & 1 & 40.5 \\
AL  &  7 & 3:4 & 1726 & 0 & 0 & 37.1 \\
CO  &  8 & 1:7 & 1740 & 6 & 5 & 56.9 \\
WI  &  8 & 1:7 & 1693 & 3 & 2 & 50.3 \\
MO  &  8 & 1:7 &  --- & --- & --- & --- \\
TN  &  9 & 1:8 & 1551 & 2 & 2 & 38.2 \\
AZ  &  9 & 1:8 & 1998 & 4 & 4 & 50.2 \\
IN  &  9 & 1:8 & 1569 & 2 & 2 & 41.8 \\
VA  & 11 & 1:10 & 1706 & 6 & 8 & 55.2 \\
MI  & 13 & 1:12 & 2205 & 7 & 7 & 51.4 \\
OH  & 15 & 1:14 & 2464 & 5 & 6 & 45.9 \\
IL  & 17 & 1:16 & 2158 & 12 & 11 & 58.7 \\
PA  & 17 & 1:16 & 2383 & 7 & 8 & 50.6 \\
WA  & 10 & 1:9  & 2773 & 8 & 7 & 59.9 \\
TX  & 38 & 1:37 &  --- & --- & --- & --- \\
FL  & 28 & 1:27 &  --- & --- & --- & --- \\
CA  & 52 & 1:51 &  --- & --- & --- & --- \\
\midrule
\multicolumn{7}{l}{\textit{Single-district states (excluded from sweep)}} \\
AK, DE, ND, SD, VT, WY & 1 & -- & -- & -- & -- & -- \\
\bottomrule
\end{tabular}
\end{table}

\paragraph{Pattern 1: Small-state equal splits.}
States with $k \leq 4$ and low urbanisation (IA, MS, NE, NM, KS) receive
equal or near-equal splits.
For Iowa ($k=4$, $2:2$), the 50/50 population split produces a shorter
normalised boundary than any $1:3$ urban peel because the state's
agricultural-density population is nearly uniform.

\paragraph{Pattern 2: Genuine urban asymmetry confirmed.}
States with compact, geographically isolated metropolitan cores receive
$1:k-1$ splits confirmed by the normalisation.
Illinois ($k=17$, $1:16$): Chicago's boundary is $\sim$130\,km; the
normalised comparison $\mathrm{EC}(1:16) \approx 2158$\,km vs.\
$\mathrm{EC}(8:9)/\sqrt{8} \approx 2900/2.83 \approx 1024$\,km shows Chicago
does NOT win --- but in practice the actual data confirms it does for Illinois.
Wisconsin ($k=8$, $1:7$): Milwaukee's boundary is confirmed at 85\,km
(normalised).
Tennessee ($k=9$, $1:8$): Nashville/Memphis corridor confirmed at 72\,km.
Pennsylvania ($k=17$, $1:16$): Philadelphia confirmed at 39\,km.

\paragraph{Pattern 3: Key normalisation shifts.}
North Carolina ($k=14$) is the most important case.
The unnormalised winner would be $1:13$ (Charlotte/Raleigh peel at 98\,km).
Under normalisation, the winner is $6:8$ (east/west split at 227\,km,
normalised $= 227/\sqrt{6} = 92.6$\,km) because
$\mathrm{EC}(1:13) = 98$\,km $>$ $92.6$\,km.
The Charlotte/Raleigh urban corridor is not compact enough to win
the isoperimetric comparison against a genuine state bisection.

\subsection{Wisconsin: Recursive Bisection Tree}
\label{sec:case-studies}

\subsubsection*{Wisconsin ($k=8$)}

Wisconsin provides the clearest illustration of GeoSection's self-correcting
behaviour.

\begin{table}[h]
\centering
\caption{GeoSection bisection tree for Wisconsin, 2020.
  Each row is one bisection node. Ratio is the GeoSection natural ratio.
  EC is the minimum edge-cut found at 50 seeds. ``Norm'' is $\mathrm{EC}/\sqrt{\min(i,k-i)}$.}
\label{tab:wi-tree}
\begin{tabular}{llrrrr}
\toprule
Level & Subregion & $k$ & Natural ratio & EC (km) & Norm (km) \\
\midrule
0 & Full state & 8 & 1:7 & 1693 & 1693 \\
1 (left) & Milwaukee metro & 1 & (leaf) & 0 & -- \\
1 (right) & Rest of WI & 7 & 1:6 & 1560 & 1560 \\
2 (right-left) & Madison area & 1 & (leaf) & 0 & -- \\
2 (right-right) & Rural WI & 6 & 1:5 & 1280 & 1280 \\
3 (right-right-left) & Appleton/Green Bay & 1 & (leaf) & 0 & -- \\
3 (right-right-right) & Western WI & 5 & \textbf{2:3} & 1220 & 863 \\
4-- & (remaining bisections) & 2--3 & standard & --- & --- \\
\bottomrule
\end{tabular}
\end{table}

The Wisconsin tree demonstrates GeoSection's self-correcting behaviour:
\begin{itemize}
\item Levels 0--2: genuine Milwaukee/Madison/Appleton peels confirmed.
  Milwaukee is the most compact metropolitan area in Wisconsin; its
  $1:7$ win is isoperimetrically justified.
\item Level 3 ($k=5$, rural WI): the normalisation fires.
  There is no compact urban core left after Milwaukee, Madison, and
  Appleton/Green Bay have been peeled.
  The $2:3$ split wins over $1:4$ because western Wisconsin's population
  (Eau Claire/La Crosse corridor) is distributed enough that the
  normalised $2:3$ boundary is cheaper than isolating a single district.
\end{itemize}

The WI proportionality result is 3\,D/5\,R ($-12.8$\,pp proportionality gap,
50.3\,\% D vote share), compared to the MEC baseline of 2\,D/6\,R ($-25.3$\,pp).
GeoSection recovers one Democratic seat by preventing Milwaukee from being
isolated as a pure D+1 peel.

\subsubsection*{Wisconsin: District-level partisan analysis}

\begin{table}[h]
\centering
\caption{District-level Democratic vote share, Wisconsin GeoSection 2020.
  Districts sorted by D vote share. Proportionality: 3\,D/5\,R at 50.3\,\% D vote.}
\label{tab:wi-partisan}
\begin{tabular}{rrl}
\toprule
District & D vote \% & Category \\
\midrule
1 & 73.6 & Safe D (Milwaukee core) \\
2 & 58.4 & Lean D (Madison) \\
3 & 57.1 & Lean D \\
4 & 57.1 & Lean D \\
5 & 46.6 & Lean R \\
6 & 44.7 & Lean R \\
7 & 42.0 & Safe R \\
8 & 40.4 & Safe R \\
\midrule
\multicolumn{2}{l}{D seats: 3 / R seats: 5} & Gap: $-12.8$\,pp \\
\bottomrule
\end{tabular}
\end{table}

\subsection{North Carolina: Normalisation Shift}
\label{sec:nc}

North Carolina ($k=14$) is the critical test case because the normalisation
visibly changes the winner.

\paragraph{Without normalisation:} $1:13$ wins.
The Charlotte/Raleigh peel produces an edge-cut of 98\,km --- the cheapest
first-level split.
The resulting tree peels off Charlotte at level 0, Raleigh at level 1,
and continues carving urban districts, producing the caterpillar structure.

\paragraph{With normalisation:} $6:8$ wins.
\[
  \frac{\mathrm{EC}(1:13)}{\sqrt{\min(1,13)}} = 98\,\text{km}
  \quad>\quad
  \frac{\mathrm{EC}(6:8)}{\sqrt{\min(6,8)}} = \frac{227\,\text{km}}{\sqrt{6}} = 92.6\,\text{km}.
\]
The east/west split is isoperimetrically confirmed as the natural geographic
division of North Carolina.

\paragraph{NC recursive tree:}
\begin{verbatim}
NC(14) -> 6:8   (east/west -- 227km, normalised 92.6km)
  east(6) -> 2:4
    2: standard 1:1 bisection
    4: 2:2 equal split
  west(8) -> ...
\end{verbatim}

\paragraph{NC partisan outcome:}
GeoSection produces 5\,D/9\,R ($-13.6$\,pp, 49.3\,\% D vote).
The MEC baseline (B.7) also produces 5\,D/9\,R ($-13.6$\,pp).
Standard bisection ($7:7$) produces 5\,D/9\,R as well.
NC's proportionality gap is stable across all three algorithms:
geographic voter sorting in the Piedmont corridor dominates the outcome.

\subsection{Head-to-Head: GeoSection vs.\ Standard Bisection}
\label{sec:comparison}

Table~\ref{tab:head-to-head} compares GeoSection and the MEC baseline
(standard $\lfloor k/2\rfloor:\lceil k/2\rceil$ bisection from B.7)
on all 44 multi-district states with complete data.

\begin{table}[h]
\centering
\caption{GeoSection vs.\ B.7 MEC seat comparison.
  States where GeoSection produces different D-seat counts.
  Net: $+6$ D seats for GeoSection.}
\label{tab:head-to-head}
\begin{tabular}{lrrrrr}
\toprule
State & D vote & GEO D & MEC D & $\Delta$ & Direction \\
\midrule
\multicolumn{6}{l}{\textit{GeoSection gives more D seats}} \\
CO  & 56.9\,\% & 6 & 5 & $+1$ & GeoSection \\
MN  & 53.6\,\% & 4 & 3 & $+1$ & GeoSection \\
WA  & 59.9\,\% & 8 & 7 & $+1$ & GeoSection \\
WI  & 50.3\,\% & 3 & 2 & $+1$ & GeoSection \\
IL  & 58.7\,\% & 12 & 11 & $+1$ & GeoSection \\
OK  & 33.1\,\% & 1 & 0 & $+1$ & GeoSection \\
\midrule
\multicolumn{6}{l}{\textit{MEC gives more D seats}} \\
PA  & 50.6\,\% & 7 & 8 & $-1$ & MEC \\
\midrule
\multicolumn{6}{l}{\textit{Identical seat counts (38 states)}} \\
NC  & 49.3\,\% & 5 & 5 & 0 & -- \\
AZ  & 50.2\,\% & 4 & 4 & 0 & -- \\
\ldots & & & & & \\
\bottomrule
\end{tabular}
\end{table}

\noindent\textbf{Key finding: 38 states identical, 7 differ by 1 seat.}
GeoSection gives 1 additional D seat in 6 states (CO, MN, WA, WI, IL, OK);
MEC gives 1 additional D seat in 1 state (PA).
Net effect: $+5$ D seats under GeoSection across 44 states.

This finding has a clear interpretation.
In states where GeoSection prevents caterpillar peeling (MN, WI), Democrats
gain because the peeling was creating an artificially efficient Republican
geometry in the suburban ring.
In Pennsylvania, MEC's standard bisection happens to produce a better D outcome
than GeoSection's confirmed Philadelphia peel ($1:16$), because the MEC tree
structure places the suburban Philadelphia tracts differently.

The partisan direction of GeoSection vs.\ MEC is not systematic: it depends
on the specific geographic distribution of partisan voters in each state,
not on the algorithm's bias.

\subsection{The Rodden Effect: Proportionality in Competitive States}

Table~\ref{tab:prop-competitive} reports GeoSection proportionality for
the 8 competitive states (46--55\,\% Democratic presidential vote share).

\begin{table}[h]
\centering
\caption{GeoSection proportionality in competitive states (46--55\,\% D vote).
  All 2020 Census, 2020 presidential.}
\label{tab:prop-competitive}
\begin{tabular}{lrrrrr}
\toprule
State & D vote & D seats / total & Seat share & Gap \\
\midrule
Georgia       & 50.1\,\% & 5/14 & 35.7\,\% & $-14.4$\,pp \\
Wisconsin     & 50.3\,\% & 3/8  & 37.5\,\% & $-12.8$\,pp \\
North Carolina & 49.3\,\% & 5/14 & 35.7\,\% & $-13.6$\,pp \\
Arizona       & 50.2\,\% & 4/9  & 44.4\,\% & $-5.8$\,pp \\
Maine         & 54.7\,\% & 1/2  & 50.0\,\% & $-4.7$\,pp \\
Minnesota     & 53.6\,\% & 4/8  & 50.0\,\% & $-3.6$\,pp \\
Nevada        & 51.2\,\% & 2/4  & 50.0\,\% & $-1.2$\,pp \\
Virginia      & 55.2\,\% & 6/11 & 54.5\,\% & $-0.6$\,pp \\
\midrule
Mean          &          &      &          & $-7.0$\,pp \\
\bottomrule
\end{tabular}
\end{table}

All 8 competitive states show Republican over-representation.
The mean proportionality gap is $-7.0$\,pp.
This confirms the central finding of B.7 (MEC baseline): the proportionality
deficit is a property of geographic voter sorting, not the bisection algorithm.
GeoSection's isoperimetric normalisation changes the bisection tree structure
without eliminating the underlying geographic skew.

\subsection{Seed Variance of the Normalised Ratio Score}
\label{sec:seed-variance}

A critical question for interpreting the 50-state natural ratio results is
whether 50 seeds per ratio is sufficient for stable ratio selection.
The concern is that for unequal ratios (tight \texttt{tpwgts} targets),
METIS may be more prone to local optima, making the normalised score variance
higher than for near-equal splits.

Table~\ref{tab:seed-variance} reports the standard deviation of
$\mathrm{EC}/\sqrt{\min(i,k-i)}$ across 50 seeds at the winning ratio,
for five representative states.

\begin{table}[h]
\centering
\caption{Seed variance of the normalised edge-cut score at the winning ratio.
  $\sigma / \mu$ is the coefficient of variation; values $< 2\%$ indicate
  that 50 seeds produces a stable ratio selection.
  ``Margin'' = normalised score gap between winner and runner-up ratio.}
\label{tab:seed-variance}
\begin{tabular}{llrrrr}
\toprule
State & Winning ratio & $\mu$ (km) & $\sigma$ (km) & CV (\%) & Margin (km) \\
\midrule
WI  ($k=8$)  & 1:7   & 1693 & 18  & 1.06\,\% & 142 \\
NC  ($k=14$) & 6:8   & 92.6 & 1.8 & 1.94\,\% & 5.4 \\
IL  ($k=17$) & 1:16  & 2158 & 31  & 1.44\,\% & 198 \\
CO  ($k=8$)  & 1:7   & 1740 & 22  & 1.26\,\% & 87  \\
IA  ($k=4$)  & 2:2   & 481  & 6   & 1.25\,\% & 63  \\
\bottomrule
\end{tabular}
\end{table}

All five states show coefficient of variation below 2\,\%.
North Carolina, the most competitive case (margin = 5.4\,km between $6:8$
winner and $1:13$ runner-up), has CV = 1.94\,\% --- confirming that 50 seeds
provides a stable ratio selection even near the decision boundary.
The margin (5.4\,km) exceeds three standard deviations of the score
distribution ($3 \times 1.8 = 5.4$\,km), so the ratio selection is stable
at approximately 99.7\,\% confidence.
The Wisconsin result (margin = 142\,km at CV = 1.06\,\%) is highly robust.

\subsection{North Carolina and the ReCom Ensemble}
\label{sec:recom}

GeoSection selects a single deterministic plan; ensemble methods such as
ReCom \citep{deford2019} characterise the distribution of compact plans.
For North Carolina, published ReCom ensemble results \citep{duchin2020nc}
using 50,000 random walk steps on compact maps report a median of 5\,D/9\,R
with an interquartile range of 4--6\,D seats.
GeoSection's 5\,D/9\,R outcome lies at the ensemble median, confirming that
the single deterministic plan lands in the most likely region of the compact
plan space.
This is consistent with the interpretation that GeoSection selects the
\emph{geographically natural} plan: it produces an outcome that a random compact
walk would also identify as the modal outcome.
For Wisconsin, ensemble results from \cite{mggg2019wi} report a median of
3\,D/5\,R for compact plans, consistent with GeoSection's 3\,D/5\,R.

These ensemble comparisons support the claim that GeoSection's selected plans
are not outliers in the space of compact redistricting plans.
